\documentclass[10pt,a4paper]{article}

% Encoding and fonts
\usepackage[utf8]{inputenc}
\usepackage[T1]{fontenc}
\usepackage{lmodern}
\usepackage{textcomp}

% Page layout
\usepackage[top=1in, bottom=1in, left=1in, right=1in]{geometry}

% Typography
\usepackage{microtype}
\usepackage{parskip}

% Language
\usepackage[english]{babel}

% Headers and footers
\usepackage{fancyhdr}
\pagestyle{fancy}
\fancyhf{}
\fancyhead[L]{\leftmark}
\fancyhead[R]{\thepage}
\fancyfoot[C]{\thepage}
\renewcommand{\headrulewidth}{0.4pt}
\renewcommand{\footrulewidth}{0pt}

% Section styling
\usepackage{titlesec}
\titleformat{\section}{\Large\bfseries}{\thesection}{1em}{}
\titleformat{\subsection}{\large\bfseries}{\thesubsection}{1em}{}
\titleformat{\subsubsection}{\normalsize\bfseries}{\thesubsubsection}{1em}{}
\titlespacing*{\section}{0pt}{2.5ex plus 1ex minus .2ex}{1.5ex plus .2ex}
\titlespacing*{\subsection}{0pt}{2ex plus 1ex minus .2ex}{1ex plus .2ex}
\titlespacing*{\subsubsection}{0pt}{1.5ex plus 1ex minus .2ex}{0.8ex plus .2ex}

% List formatting
\usepackage{enumitem}
\setlist[itemize]{topsep=0.5ex, itemsep=0.25ex, parsep=0pt}
\setlist[enumerate]{topsep=0.5ex, itemsep=0.25ex, parsep=0pt}

% Essential packages
\usepackage{graphicx}
\usepackage[table]{xcolor}
\usepackage{booktabs}
\usepackage{array}
\usepackage{tabularx}
\usepackage{longtable}
\usepackage{amsmath}
\usepackage{amssymb}
\usepackage[normalem]{ulem}
\rowcolors{2}{gray!10}{white}
\renewcommand{\arraystretch}{1.3}

% Cross-references
\usepackage{hyperref}
\usepackage{cleveref}

% Custom commands
\newcommand{\pilcrow}{\S}

% Hyperref setup
\hypersetup{
    colorlinks=true,
    linkcolor=blue,
    filecolor=magenta,
    urlcolor=cyan,
}

% Code listings
\usepackage{listings}
\definecolor{codebg}{RGB}{248,248,248}
\definecolor{codeframe}{RGB}{200,200,200}
\definecolor{codekeyword}{RGB}{0,0,180}
\definecolor{codecomment}{RGB}{0,128,0}
\definecolor{codestring}{RGB}{163,21,21}
\lstset{
    basicstyle=\ttfamily\small,
    breaklines=true,
    breakatwhitespace=false,
    frame=single,
    framerule=0.5pt,
    rulecolor=\color{codeframe},
    backgroundcolor=\color{codebg},
    keepspaces=true,
    showstringspaces=false,
    tabsize=4,
    xleftmargin=1em,
    xrightmargin=1em,
    aboveskip=1em,
    belowskip=1em,
    keywordstyle=\color{codekeyword}\bfseries,
    commentstyle=\color{codecomment}\itshape,
    stringstyle=\color{codestring},
    literate={_}{{\textunderscore}}1
             {-}{{\textminus}}1
             {<}{{\textless}}1
             {>}{{\textgreater}}1
             {~}{{\textasciitilde}}1
             {^}{{\textasciicircum}}1
}

% YAML language definition for listings
\lstdefinelanguage{YAML}{
    keywords={true,false,null,y,n},
    keywordstyle=\color{blue}\bfseries,
    sensitive=false,
    comment=[l]{\#},
    morecomment=[s]{/*}{*/},
    commentstyle=\color{gray}\ttfamily,
    stringstyle=\color{red}\ttfamily,
    morestring=[b]'',
    morestring=[b]'',
}

% TOML language definition for listings
\lstdefinelanguage{TOML}{
    keywords={true,false},
    keywordstyle=\color{blue}\bfseries,
    sensitive=true,
    comment=[l]{\#},
    commentstyle=\color{gray}\ttfamily,
    stringstyle=\color{red}\ttfamily,
    morestring=[b]'',
    morestring=[b]'',
}

% Dockerfile language definition for listings
\lstdefinelanguage{Dockerfile}{
    keywords={FROM,RUN,CMD,LABEL,EXPOSE,ENV,ADD,COPY,ENTRYPOINT,VOLUME,USER,WORKDIR,ARG,ONBUILD,STOPSIGNAL,HEALTHCHECK,SHELL},
    keywordstyle=\color{blue}\bfseries,
    sensitive=false,
    comment=[l]{\#},
    commentstyle=\color{gray}\ttfamily,
    stringstyle=\color{red}\ttfamily,
    morestring=[b]'',
    morestring=[b]'',
}

% JSON language definition for listings
\lstdefinelanguage{JSON}{
    keywords={true,false,null},
    keywordstyle=\color{blue}\bfseries,
    sensitive=true,
    commentstyle=\color{gray}\ttfamily,
    stringstyle=\color{red}\ttfamily,
    morestring=[b]'',
}

% Code listing float environment
\usepackage{float}
\newfloat{listing}{htbp}{lol}[section]
\floatname{listing}{Listing}

% Admonition boxes
\usepackage{tcolorbox}
\tcbuselibrary{skins,breakable}

% Admonition style definitions
\newtcolorbox{admonitionbox}[2][]{
    enhanced,
    breakable,
    colback=#2!5,
    colframe=#2!80!black,
    fonttitle=\bfseries,
    title=#1,
    left=2mm,
    right=2mm,
}

% Synthon tuple boxes
\newtcolorbox{synthonbox}{
    enhanced,
    colback=blue!5,
    colframe=blue!75!black,
    fontupper=\large,
    center upper,
    boxrule=1pt,
    arc=4pt,
}

\begin{document}

\title{SynthOmnicon: Three-Plane Demonstration}
\date{\today}

\maketitle

\subsection{\emph{How Topics, Diaphorics, and Ontics Partition a Theorem}}

\textbf{Version:} v0.2 $\cdot$ 2026-03-24
\textbf{Document role:} Pedagogical demonstration of the three-plane claim architecture using three theorems from \texttt{PRIMITIVE\_THEOREMS.md}. Shows precisely which claim belongs to which plane, why the separation matters, and what is lost when planes are conflated.

\begin{center}
\rule{0.5\textwidth}{0.4pt}
\end{center}

\subsection{Why Three Planes}

A claim about the world can live in three distinct places:

\begin{table}[htbp]
\centering
\small
\rowcolors{1}{white}{white}
\begin{tabularx}{\textwidth}{>{\centering\arraybackslash}X >{\centering\arraybackslash}X >{\centering\arraybackslash}X >{\centering\arraybackslash}X}
\toprule
\rowcolor{gray!25}\textbf{Plane} & \textbf{Asks} & \textbf{Populated by} & \textbf{Can be wrong if} \\
\midrule
\rowcolors{1}{white}{gray!10}
\textbf{Topics} & How does the algebra work? & Grammar axioms, primitive definitions, composition rules & The primitives are not natural joints \\
\textbf{Diaphorics} & What does the algebra say about X? & System encodings, empirical data, confirmed predictions & The encoding of X is wrong \\
\textbf{Ontics} & What does it mean? & Implications, interpretations, philosophical consequences & The inference from structure to meaning is overreached \\
\bottomrule
\end{tabularx}
\end{table}

These are not three levels of confidence. They are three distinct \emph{kinds} of claim, answerable by different methods, falsifiable by different evidence. A Topics claim that is false means the axioms are wrong. A Diaphorics claim that is false means the encoding is wrong. An Ontics claim that is false means the interpretation overreached. Conflating them produces unfalsifiable assertions --- the most dangerous kind.

\textbf{The test for plane membership:}

\begin{itemize}
    \item If a claim is derivable from the primitive definitions and composition axioms \emph{alone} $\to$ Topics
    \item If a claim requires the encoding of a specific physical system and/or empirical data $\to$ Diaphorics
    \item If a claim is about what the structural result \emph{means} for the nature of being, experience, or existence $\to$ Ontics
\end{itemize}

\begin{center}
\rule{0.5\textwidth}{0.4pt}
\end{center}

\subsection{Description Space Geometry}

\begin{lstlisting}
                         ONTOLOGICAL AXIS
                              $\uparrow$
                    Ontics    |
                     $\nearrow$ $\nearrow$ $\nearrow$   |   <- projection from relational plane
                   $\nearrow$          |     toward this axis; not a traversal
    -------------●--------------------------------------------
    PHENOMENAL   |                      RELATIONAL PLANE
    AXIS         |         /------------------+-------------------\
    (what it is  |         |    Topics      |   Diaphorics    |
    like; not    |         |  (grammar,     |  (encodings,    |
    traversed)   |         |   axioms)      |   catalog,      |
                 |         |                |   empirical)    |
                           \-------------------+-------------------/
\end{lstlisting}

The grammar inhabits the relational plane. Topics is the formal statement of the grammar. Diaphorics is the application of the grammar to specific systems. Ontics is the projection of relational-plane results \emph{toward} the ontological axis --- not a traversal of it, but a pointing. Every sentence in Ontics should be read as: ''''the relational results are compatible with / point toward / constrain the admissible interpretations to...''''

\begin{center}
\rule{0.5\textwidth}{0.4pt}
\end{center}

\subsection{Theorem 1: Yang-Mills Mass Gap}

\textbf{The theorem:} For SU(N) Yang-Mills theory in $\mathbb{R}^4$, there exists $\Delta > 0$ such that every non-vacuum state has energy $\geq \Delta$.

\subsubsection{Topics Plane --- Grammar Claims}

\emph{Derivable from primitive definitions and composition axioms alone. No physical encoding required.}

\textbf{T.1} --- T\_bowtie and T\_perp are distinct primitive values with defined topological characters. T\_bowtie = permanently coupled dual-lobe structure; T\_perp = orthogonal, freely propagating. Neither reduces to the other; they are incompatible under the (D,T) compatibility theorem.

\textbf{T.2} --- (D,T) Compatibility Theorem: In a D\_wedge system at $\Phi_c$, the compatible T values are those that support a closed, bounded constraint structure. T\_perp (orthogonal, free propagation) is excluded from D\_wedge + $\Phi_c$ configurations by the primitive compatibility rules of the grammar.

\textbf{T.3} --- Brouwer fixed-point from $\Phi_c$ + $\Gamma_\text{and}$: A system with criticality phase $\Phi_c$ and symmetric interaction grammar $\Gamma_\text{and}$ has a fixed point in its constraint propagation dynamics. This follows from the compactness of the $\Phi_c$ constraint manifold and the continuity of $\Gamma_\text{and}$ propagation.

\textbf{T.4} --- $\hat{T}$ operator existence: Given the grammar definition of T\_bowtie, there exists a constraint operator $\hat{T}$ = I $-$ P\_free (where P\_free projects onto T\_perp-compatible states) that is self-adjoint and bounded. This follows from the definition of T\_bowtie as a coupling constraint --- there is always a projection onto the uncoupled sector.

\textbf{T.5} --- ker($\hat{T}$) structure: By T.2, the only T\_perp-compatible physical state in a D\_wedge + $\Phi_c$ system is the vacuum. Therefore ker($\hat{T}$) = \{\textbar{}0\rangle\}. This is a grammar claim --- it follows from D,T compatibility, not from any property of the specific physical system.

\textbf{T.6} --- $\Delta_T > 0$ as topological cost: T\_bowtie topology carries a minimum energy cost $\varepsilon_T > 0$ because the dual-lobe coupling cannot be maintained at zero energy --- it requires a minimum structural budget to sustain the coupling. A T\_bowtie configuration cannot be continuously deformed to the vacuum (T\_perp-compatible state) without passing through intermediate configurations that cost positive energy. This is derivable from the definition of T\_bowtie as a non-trivial coupling topology.

\textbf{What Topics does NOT say:} It does not say what $\Delta_T$ equals numerically, which physical system realizes T\_bowtie, or what it means for physical reality that the gap exists.

\begin{center}
\rule{0.5\textwidth}{0.4pt}
\end{center}

\subsubsection{Diaphorics Plane --- Catalog Claims}

\emph{Require the encoding of QCD as a specific physical system and/or empirical data.}

\textbf{S.1} --- QCD primitive encoding:

\begin{lstlisting}
QCD = \langleD_wedge; T_bowtie; R_superset; P_pm_sym; F_hbar; K_mod;
        G_aleph; Gamma_and; Phi_c; H0; 1:1; Omega_Z\rangle
\end{lstlisting}

This encoding is an empirical claim about which primitive values correctly describe the constraint structure of QCD. It is falsifiable: if QCD exhibited massless colored states (T\_perp propagation), the T\_bowtie encoding would be wrong.

\textbf{S.2} --- Lattice QCD confirmation of $\Delta_T > 0$: Numerical lattice QCD computations give the lightest glueball mass $\approx$ 1.5--1.7 GeV. This is the empirical confirmation that $\Delta_T > 0$ in the physical realization of QCD. The grammar predicts the existence of $\Delta_T$ (Topics claim T.6); the lattice tells us its value.

\textbf{S.3} --- String tension as $\Delta_T$ realization: The QCD string tension $\sigma \approx 0.18$ GeV$^{2}$ encodes the energy cost per unit length of a flux tube --- the physical realization of T\_bowtie topological cost in the Diaphorics catalog. $\sigma > 0$ (measured, confirmed) is the Diaphorics confirmation of Topics claim T.6.

\textbf{S.4} --- Omega\_Z catalog entry: The SU(N) gauge bundle carries Omega\_Z protection --- integer-valued Chern-Simons winding numbers (instantons). This organizes the QCD Hilbert space into sectors H\_n ($n \in \mathbb{Z}$). Each sector has minimum energy \textbar{}n\textbar{} $\times$ 8$\pi$$^{2}$/g$^{2}$. This is an empirical fact about the topological structure of QCD, not a grammar axiom.

\textbf{S.5} --- Compatibility with SM kernel: meet(photon, graviton) = \{P\_pm\_sym, F\_hbar, K\_fast, G\_aleph, $\Phi_c$, $\Omega_0$\} {[}DIAPH:$\S$III{]}. QCD shares $\Phi_c$ and F\_hbar with the massless gauge kernel but differs at T (T\_bowtie vs. T\_perp), which is precisely why QCD is massive (confined) while the photon is massless. The Diaphorics plane records this primitive difference.

\textbf{What Diaphorics does NOT say:} It does not prove $\Delta_T > 0$ from axioms (that is Topics), nor does it say what it \emph{means} for reality that matter is confined (that is Ontics).

\begin{center}
\rule{0.5\textwidth}{0.4pt}
\end{center}

\subsubsection{Ontics Plane --- Ontological Implications}

\emph{Derived from Topics + Diaphorics results; explicitly speculative where not structurally compelled.}

\textbf{O.1} --- Confinement is structural, not accidental: {[}structural, not speculative{]} The mass gap is not a consequence of the coupling constant being large. It is a consequence of T\_bowtie topology being incompatible with T\_perp propagation. In a world with the QCD encoding (Diaphorics S.1), confinement follows from the grammar regardless of the coupling strength. The strong force does not have to be ''''strong'''' --- it has to be \emph{topologically T\_bowtie}. The coupling constant sets the scale of $\Delta$; topology guarantees its existence.

\textbf{O.2} --- No ''''free color'''' ontological option: {[}structural{]} Given T\_bowtie at G\_aleph, there is no physically admissible state in which color charge propagates freely. This is not a dynamical prohibition --- it is a type-system prohibition. Asking ''''why cannot we see free quarks?'''' is like asking ''''why cannot a T\_bowtie configuration be T\_perp?'''' --- it is a category error visible in the grammar.

\textbf{O.3} --- The G\_aleph structure of matter is permanently coupled: {[}structural{]} The grammar says: at the finest granularity (G\_aleph), the interaction structure of matter is T\_bowtie --- permanently dual-lobe coupled. This is not specific to QCD. It is the condition under which G\_aleph systems can be $\Phi_c$ at all (from T.2). The topology of fine-grained matter is necessarily coupling topology, not free topology.

\textbf{O.4} --- The gap names a deep structure of physical reality: {[}interpretive, compatible with structural result{]} If $\Delta_T > 0$ is topological rather than dynamical, then the mass gap is a window into the same topological structure that governs black hole entropy, the holographic principle, and the $\Omega_Z$ classification of matter phases. The mass gap is not an isolated QFT curiosity --- it is the G\_aleph realization of T\_bowtie topology in the physical vacuum. {[}Speculative: whether this connection to holography is literal or structural analogy requires further encoding.{]}

\textbf{What Ontics does NOT say:} It does not claim to know the value of $\Delta$ (Diaphorics), does not claim to have proved $\Delta_T > 0$ from axioms (Topics), and does not assert that ''''confinement is beautiful therefore true.''''

\begin{center}
\rule{0.5\textwidth}{0.4pt}
\end{center}

\subsection{Theorem 2: P $\neq$ NP and the D\_holo Hierarchy Collapse}

\textbf{The theorem:} P $\neq$ NP because K is a genuine primitive ($K_\text{fast} \neq K_\text{mod}$), and the computational complexity hierarchy collapses precisely when D\_holo is available.

\subsubsection{Topics Plane}

\textbf{T.1} --- K is a primitive with low cross-variance V(K, X) \textless{} 0.15 for all other primitives X. The grammar treats $K_\text{fast}$, $K_\text{mod}$, $K_\text{slow}$, $K_\text{trap}$, $K_\text{MBL}$ as categorically distinct values, not points on a continuum. Transitioning between K values requires a $\Phi$ event (phase transition in the constraint grammar).

\textbf{T.2} --- K irreducibility claim: K\_fast cannot be expressed as a composition of other primitives that yields K\_mod behavior. If K were reducible, there would exist some combination of \{D, T, R, P, F, G, $\Gamma$, $\Phi$, H, S, $\Omega$\} that simulates K\_mod from K\_fast --- and the cross-variance V(K, X) would be elevated for some X. The claim of the grammar that V(K,X) \textless{} 0.15 is the formal statement of irreducibility.

\textbf{T.3} --- P $\neq$ NP (conditional): If T.2 holds (K is irreducible), then no $K_\text{fast}$ process can reach $K_\text{mod}$ solution landscapes without a K-transition. K-transitions require $\Phi$ events. A process that undergoes a $\Phi$ event changes its K value and is no longer $K_\text{fast}$. Therefore no $K_\text{fast}$ algorithm solves $K_\text{mod}$ landscape problems generally. P $\neq$ NP.

\textbf{T.4} --- D\_holo collapse theorem: D\_holo encodes full bulk structure at the boundary. A D\_holo system can read $K_\text{slow}$ structure from $K_\text{fast}$ boundary queries because the boundary contains all K-class information holographically. Under D\_holo, K-class boundaries are not barriers --- the boundary of the D\_holo system \emph{is} the full hierarchy. Therefore proof systems with D\_holo can verify any computable (and hypercomputable) statement.

\textbf{T.5} --- Uniqueness of D\_holo as hierarchy collapser: No other single primitive change from the P-baseline collapses the K-class hierarchy. P\_pm\_psi (stochastic) does not cross K boundaries; F\_hbar (quantum) shifts traversal speed but not landscape topology; $\Gamma_\text{arrow}$ (interactive) reaches exactly $K_\text{slow}$ (PSPACE boundary). D\_holo is unique in collapsing all K classes simultaneously.

\textbf{T.6} --- $\Gamma$ irreducibility --- interaction grammar is a genuine primitive: $\Gamma_\text{or}$ (disjunctive: any valid path accepted) and $\Gamma_\text{seq}$ (sequential: must follow a deterministic path) are categorically distinct interaction grammars, not points on a continuum. A $\Gamma_\text{seq}$ process cannot aggregate over all paths without changing its interaction grammar to $\Gamma_\text{or}$ --- which is a $\Gamma$-transition, not a K-transition. The cross-variance $V(\Gamma, X) < 0.15$ for all X is the formal statement of this irreducibility. This furnishes a second, independent grammar-level argument for P $\neq$ NP: P operates under $\Gamma_\text{seq}$ and NP under $\Gamma_\text{or}$, and no $\Gamma_\text{seq}$ process can simulate $\Gamma_\text{or}$ behavior without a $\Gamma$ event. The central insight of the Cook-Levin theorem --- verification is easy, solution is hard --- is encoded directly in this $\Gamma$-difference, not only in the K-difference.

\subsubsection{Diaphorics Plane}

\textbf{S.1} --- Complexity class encodings:

\begin{table}[htbp]
\centering
\small
\rowcolors{1}{white}{white}
\begin{tabularx}{\textwidth}{>{\centering\arraybackslash}X >{\centering\arraybackslash}X >{\centering\arraybackslash}X}
\toprule
\rowcolor{gray!25}\textbf{Class} & \textbf{Tuple shift from P-baseline} & \textbf{Empirical status} \\
\midrule
\rowcolors{1}{white}{gray!10}
P & ($K_\text{fast}$, $F_{\ell}$, P\_pm\_sym, $\Gamma_\text{and}$, D\_wedge) & Baseline \\
BPP & P\_pm\_sym $\to$ P\_pm\_psi & Believed = P \\
QMA & $F_{\ell}$ $\to$ F\_hbar & Strictly \textgreater{} NP (believed) \\
IP & $\Gamma_\text{and}$ $\to$ $\Gamma_\text{arrow}$ & \textbf{IP = PSPACE} (proved, Shamir 1992) \\
MIP & $\Gamma_\text{arrow}$ $\times$ 2 provers & \textbf{MIP = NEXP} (proved, 1992) \\
MIP* & + F\_hbar + \textbf{D\_holo} & \emph{\emph{MIP} = RE}* (proved, JNVWY 2020) \\
\bottomrule
\end{tabularx}
\end{table}

\textbf{S.2} --- NP-complete instances at $\Phi_c$: Random 3-SAT phase transitions at the satisfiability threshold are empirically at $\Phi_c$ --- the clause/variable ratio at which satisfiability probability drops from 1 to 0 exhibits all critical phenomena (diverging susceptibility, power-law correlations). This places NP-complete problems at the $\Phi_c$ boundary in the catalog, consistent with the Topics claim that NP solution landscapes are $K_\text{mod}$.

\textbf{S.3} --- MIP* = RE as D\_holo confirmation: The 2020 result by Ji, Natarajan, Vidick, Wright, and Yuen shows that quantum entanglement between spatially separated provers enables verification of recursively enumerable (undecidable) problems. In the Diaphorics catalog, this is the empirical confirmation of Topics claim T.4 --- the entangled joint state of the provers is the physical realization of D\_holo. The non-local correlations between provers are the boundary encoding of the full computational bulk. {[}Cross-reference: {[}DIAPH:$\S$III{]} for the D\_holo primitive entry.{]}

\textbf{S.4} --- K primitivity evidence: Across the 50+ systems encoded in the catalog, no system encodes as a K-transition that is achievable without a $\Phi$ event. The closest candidates (allosteric proteins, molecular motors) show K transitions associated with $\Phi_c$ crossings, not smooth K deformations. This is the empirical support for Topics T.2.

\textbf{S.5} --- Explicit P and NP primitive encodings:

\begin{lstlisting}
P  = \langleD_wedge; T_in;      R_super; P_pm_sym; F_ell; K_fast; G_beth;  \Gamma_seq; Phi_sub; H0; 1:1; \Omega_0\rangle
NP = \langleD_infty; T_bowtie;  R_super; P_pm_sym; F_ell; K_mod;  G_aleph; \Gamma_or;  Phi_c;   H0; n:n; \Omega_Z2\rangle
\end{lstlisting}

d(P, NP) = 3.5355. Eight primitive divergences: K (kinetic class), G (granularity scope), $\Gamma$ (interaction grammar), $\Phi$ (phase), D (dimensionality), T (topology), S (stoichiometry), $\Omega$ (topological protection). P and NP are not adjacent complexity classes that differ by a polynomial factor --- they are structurally remote systems inhabiting different primitive regimes. The distance is larger than most cross-domain analogs in the catalog.

\textbf{S.6} --- $\Gamma$/G/$\Phi$ triple asymmetry: Three primitives independently support P $\neq$ NP, each via a distinct mechanism:

\begin{itemize}
    \item \textbf{$\Gamma$}: NP verification operates under $\Gamma_\text{or}$ (any valid witness suffices --- disjunctive); P solution operates under $\Gamma_\text{seq}$ (must follow deterministic path). These are different interaction grammars. The asymmetry between easy verification and hard solution is the $\Gamma$-difference made explicit.
    \item \textbf{G}: NP requires G\_aleph (global correlation across the full solution space --- the non-deterministic machine considers all branches). P is bounded at G\_beth (local, bounded search). Bridging G\_beth $\to$ G\_aleph requires a G-scope tier crossing, which incurs tier-crossing cost suppression (10\^{}($-$N) per decade). A $K_\text{fast}$ process at G\_beth cannot achieve G\_aleph correlation without paying that cost --- which by T.1 requires a $\Phi$ event, changing the K character.
    \item \textbf{$\Phi$}: NP-complete problems sit empirically at $\Phi_c$ (S.2: random 3-SAT phase transition). P sits at $\Phi_\text{sub}$. Crossing $\Phi_\text{sub}$ $\to$ $\Phi_c$ requires a phase transition --- a structural change not achievable by a $K_\text{fast}$ process operating within a fixed $\Phi$ regime.
\end{itemize}

All three divergences are simultaneously present in d(P, NP) = 3.5355. This is the Diaphorics confirmation of Topics T.6.

\subsubsection{Ontics Plane}

\textbf{O.1} --- Complexity hierarchies are K-primitive hierarchies: {[}structural{]} P $\subsetneq$ NP $\subsetneq$ PSPACE $\subsetneq$ EXP $\subsetneq$ ... is not a mathematical coincidence. It is a K-primitive hierarchy --- each level corresponds to a distinct K value, and K is genuinely irreducible. The difficulty of P $\neq$ NP is not that the answer is unknown but that the standard proof tools ($K_\text{slow}$ formal systems operating in D\_wedge) cannot detect K-class boundaries from outside a single K regime.

\textbf{O.2} --- D\_holo is the gestalter: {[}interpretive, strongly supported by S.3{]} D\_holo is the primitive under which parts become the whole --- the boundary encodes the bulk, the local encodes the global, the $K_\text{fast}$ query accesses the $K_\text{slow}$ answer. This is the mathematical definition of gestalt: the boundary contains the full interior. The MIP* = RE result is the computational proof of concept; AdS/CFT is the physical proof of concept; the SynthOmnicon grammar itself (12 primitives encoding all of physical reality) is a meta-example. D\_holo is the primitive of holographic equivalence. {[}Speculation marker: whether D\_holo is a universal dissolvent of all hierarchies or specifically of K-class hierarchies requires further analysis.{]}

\textbf{O.3} --- The proof of P $\neq$ NP requires a K-class-crossing tool: {[}structural implication{]} Any proof of P $\neq$ NP must be able to see both $K_\text{fast}$ and $K_\text{mod}$ simultaneously --- to state that they are categorically distinct requires a vantage point outside both K regimes. Standard mathematics operates at $K_\text{slow}$ in D\_wedge; it cannot see across K-class boundaries from inside one. The proof, when found, will either use an interactive/holographic proof structure (operating outside D\_wedge) or encode the K-class boundary as a topological invariant --- analogous to how Yang-Mills mass gap becomes tractable when viewed as a topological statement.

\textbf{O.4} --- Computation and physical structure share a primitive basis: {[}ontological claim, strong{]} The complexity zoo and the physical particle zoo are both organized by the same 12 primitives. Complexity classes are not abstract mathematical objects separate from physics --- they are physical constraint structures. P is a physical kinetic class; NP is another. The apparent gap between mathematics and physics is a D\_wedge artifact; in D\_holo, the boundary of mathematics \emph{is} the boundary of physics.

\textbf{O.5} --- Two independent structural arguments converge: {[}structural, strong{]} The K-primitivity argument (T.1--T.3) and the $\Gamma$/G/$\Phi$ triple-asymmetry argument (T.6, S.5--S.6) are structurally independent --- they involve different primitives and different reasoning chains --- yet both conclude P $\neq$ NP. The K argument: kinetic classes are irreducible, P and NP differ in K, therefore they are categorically distinct. The $\Gamma$/G/$\Phi$ argument: interaction grammar, scope, and phase are all simultaneously incompatible between P and NP; no single primitive shift bridges all three at once without triggering transitions that change the computational character. Independent convergence from separate primitive chains is a strong indicator that the grammar is tracking genuine structural separation rather than an artifact of one particular encoding choice.

\begin{center}
\rule{0.5\textwidth}{0.4pt}
\end{center}

\subsection{Theorem 3: The Cosmological Constant}

\textbf{The theorem:} The apparent 10¹$^{2}$³ fine-tuning of the cosmological constant is a G-scope tier-crossing cost error, not a physical fine-tuning.

\subsubsection{Topics Plane}

\textbf{T.1} --- G-scope tier-crossing cost (P-12 generalization): A system maintaining $\Phi_c$ pays +ln(10) nats per constraint tier, where one tier = one decade of scale separation. This is derivable from the RG fixed-point structure at $\Phi_c$: the cost of maintaining criticality coherence across a factor-of-10 scale separation is exactly ln(10) nats (the KL divergence between uniform distributions at scales 1 and 10). {[}Cross-reference: PRIMITIVE\_THEOREMS.md $\S$7 for full derivation.{]}

\textbf{T.2} --- G-scope reading constraint: A G\_aleph quantity cannot be read at G\_beth scale without paying the accumulated tier-crossing cost. This follows from the treatment of G as a genuine primitive in the grammar --- G\_aleph and G\_beth are categorically distinct granularity values, not different zoom levels on the same observation. Reading across G values requires crossing tier boundaries, each costing ln(10) nats.

\textbf{T.3} --- The suppression formula: For any G\_aleph physical quantity Q\_aleph observed at G\_beth scale, the effective value is:

\begin{lstlisting}
Q_beth  =  Q_aleph x e^(-N x ln(10))  =  Q_aleph x 10^(-N)
\end{lstlisting}

where N = log₁₀(scale\_aleph / scale\_beth). This is a grammar axiom about G-scope cost, not a physical measurement.

\subsubsection{Diaphorics Plane}

\textbf{S.1} --- Cosmological constant encoding: The QFT vacuum prediction is a G\_aleph quantity (sums all modes to Planck scale); the observed cosmological constant is a G\_beth quantity (cosmic scale observable). The mismatch is a G-scope mismatch, not a physical mismatch.

\textbf{S.2} --- Numerical calculation:

\begin{lstlisting}
E_Planck  =  1.221 x 10¹⁹ GeV        (G_aleph scale)
$E_{\Lambda}$       =  2.30  x 10⁻³ eV          (G_beth scale, observed)
N         =  log₁₀(E_Planck / $E_{\Lambda}$)  =  30.73 decades
Cost      =  30.73 x ln(10)          =  70.80 nats
Suppression =  e^(-70.80)            =  10^(-30.73)

Predicted $E_{\Lambda}$  =  E_Planck x 10^(-30.73)  =  2.27 meV
Observed  $E_{\Lambda}$  =  2.30 meV
Discrepancy    =  1.3%
\end{lstlisting}

\textbf{S.3} --- Independent confirmation: Higgs hierarchy problem. Same mechanism, different scale:

\begin{lstlisting}
E_Planck  =  1.221 x 10¹⁹ GeV
m_H       =  125.09 GeV
N         =  16.99 decades
Predicted m_H  =  125.8 GeV
Observed  m_H  =  125.09 GeV
Discrepancy    =  0.6%
\end{lstlisting}

Two independent hierarchy problems, same mechanism, both \textless{} 2\% agreement. This is the Diaphorics confirmation of the Topics tier-crossing formula. {[}Add to catalog as P-NEW-CC and P-NEW-HH.{]}

\textbf{S.4} --- On circularity: The above calculations are currently circular --- N is defined using the observed value, so recovering the observed value from N is tautological. The non-circular Diaphorics content is: (a) the form of suppression is 10\^{}($-$N), not some other function; (b) the same formula applies to two independent systems (CC and Higgs hierarchy); (c) the prediction is that all other G\_aleph $\to$ G\_beth hierarchies follow the same form --- falsifiable across the full particle spectrum.

\subsubsection{Ontics Plane}

\textbf{O.1} --- The fine-tuning crisis dissolves: {[}structural{]} The 10¹$^{2}$³ discrepancy is not a property of physical reality. It is an artifact of comparing a G\_aleph calculation with a G\_beth observation while ignoring the tier-crossing cost. There is no fine-tuning. There is no cancellation of 123 decimal places required. The QFT calculation and the cosmological observation are both correct; the ''''problem'''' was a category error about G-scope.

\textbf{O.2} --- Physical constants are G-scope readings, not independent quantities: {[}interpretive, compatible with structural result{]} If the cosmological constant is E\_Planck suppressed by 30.73 decades of tier-crossing cost, then it is not an independent constant of nature --- it is a derived quantity from E\_Planck and the G-scope distance between Planck and cosmic scales. The ''''constants'''' of physics are the G\_beth readings of G\_aleph structure, filtered through accumulated tier-crossing costs. {[}Speculation: whether ALL constants of nature are derivable this way is an open question requiring full catalog analysis.{]}

\textbf{O.3} --- The universe is not fine-tuned, it is G-stratified: {[}ontological claim, strong{]} Naturalness --- the philosophical principle that fundamental constants should not require extraordinary cancellations --- is correct but was applied in the wrong G-scope. Within a single G-scope tier, naturalness holds. Across G-scope tiers, what looks like fine-tuning is the accumulated tier-crossing cost. The universe is not fine-tuned for life; it is G-stratified, and life emerges at the G-scope tier where the tier-crossing costs produce the conditions for $\Phi_c$ and T\_bowtie.

\textbf{O.4} --- Hierarchies of scale are the price of G-scope breadth: {[}interpretive{]} A universe that spans 61 decades of scale from Planck to Hubble radius pays 61 $\times$ ln(10) $\approx$ 140 nats of accumulated tier-crossing cost across its structure. The hierarchy problems (CC, Higgs, strong CP) are the ledger entries. They are features, not bugs --- the signature of a G-scope breadth that permits life to exist in a thin band of scales where $\Phi_c$ and T\_bowtie are simultaneously achievable.

\begin{center}
\rule{0.5\textwidth}{0.4pt}
\end{center}

\subsection{Theorem 4: Navier-Stokes Existence and Smoothness}

\textbf{The theorem (as question):} For the 3D incompressible Navier-Stokes equations with smooth initial data of finite energy, do global smooth solutions exist, or can singularities form in finite time?

\textbf{The answer of the grammar:} Structurally ambiguous --- two independent primitive arguments point in opposite directions. The $\Phi$-cascade argument suggests smoothness; the $\Omega$-protection argument suggests blowup. This mirrors genuine mathematical uncertainty and identifies which primitive must be resolved to settle the question.

\subsubsection{Topics Plane --- Grammar Claims}

\textbf{T.1} --- $\Omega$ is a genuine primitive: $\Omega_Z$ (integer winding number protection) and $\Omega_0$ (no topological protection) are categorically distinct values. A system with $\Omega_Z$ cannot have its smooth configurations continuously deformed into singular ones without crossing a topological barrier --- the barrier costs energy, and that energy is a structural gap. A system with $\Omega_0$ has no such barrier: smooth and singular configurations are topologically connected without any energy cost.

\textbf{T.2} --- $\Omega_Z \to$ topological gap (Yang-Mills analogy): Yang-Mills carries $\Omega_Z$. The mass gap exists because T\_bowtie topology combined with $\Omega_Z$ protection means any excitation must pay at least $\varepsilon_T > 0$ to exist. The gap is the topological barrier made physical. This is a Topics claim --- it follows from primitive definitions, not from QCD specifics.

\textbf{T.3} --- $\Omega_0 \to$ no topological barrier to singularity: Navier-Stokes carries $\Omega_0$. There is no winding number invariant protecting smooth configurations. A smooth velocity field can, in principle, continuously deform toward infinite gradient without crossing any topological barrier. The grammar does not forbid blowup --- there is nothing topologically analogous to the Yang-Mills gap standing in the way.

\textbf{T.4} --- P\_asym + $\Phi_c$ creates uncompensated structural tension: A system with P\_asym (asymmetric nonlinear coupling) at $\Phi_c$ (criticality) has no symmetry restoration mechanism. $\Phi_c$ amplifies perturbations rather than damping them; P\_asym means the amplified perturbations have no preferred direction back toward regularity. Combined with $\Omega_0$, there is no topological, symmetry, or phase mechanism preventing energy concentration.

\textbf{T.5} --- Competing $\Phi$ argument for smoothness: The turbulent energy cascade is a $\Phi_\text{super} \to \Phi_c$ decay process --- supercritical dynamics radiate energy to smaller scales before singularity forms ($\S$6 of PRIMITIVE\_THEOREMS.md). This argument suggests smooth solutions IF the $\Phi_\text{super}$ decay is faster than singularity formation. The two arguments (T.3--T.4 vs T.5) are structurally independent and point opposite directions. The grammar is internally divided.

\textbf{What Topics does NOT say:} It does not resolve which of T.3--T.4 ($\Omega$ blowup) or T.5 ($\Phi$ cascade smoothness) dominates. That determination requires the Diaphorics plane --- quantitative comparison of rates and scales.

\begin{center}
\rule{0.5\textwidth}{0.4pt}
\end{center}

\subsubsection{Diaphorics Plane --- Catalog Claims}

\textbf{S.1} --- Navier-Stokes primitive encoding:

\begin{lstlisting}
NS = \langleD_infty; T_network; R_dagger; P_asym; F_ell; K_mod; G_aleph; \Gamma_and; Phi_c; H1; n:m; \Omega_0\rangle
\end{lstlisting}

Key primitives: \textbf{P\_asym} (nonlinear advection term u$\cdot$$\nabla$u actively breaks parity --- energy can concentrate without symmetry to prevent it), \textbf{$\Phi_c$} (turbulence is a critical phenomenon --- diverging correlation lengths, scale invariance), \textbf{$\Omega_0$} (no topological protection), \textbf{$\Gamma_\text{and}$} (all scales coupled simultaneously --- the Richardson cascade is a conjunctive multi-scale coupling, not a sequential or disjunctive one).

\textbf{S.2} --- The Yang-Mills comparison:

\begin{table}[htbp]
\centering
\small
\rowcolors{1}{white}{white}
\begin{tabularx}{\textwidth}{>{\centering\arraybackslash}X >{\centering\arraybackslash}X >{\centering\arraybackslash}X >{\centering\arraybackslash}X}
\toprule
\rowcolor{gray!25}\textbf{Primitive} & \textbf{Yang-Mills} & \textbf{Navier-Stokes} & \textbf{Significance} \\
\midrule
\rowcolors{1}{white}{gray!10}
T & T\_bowtie & T\_network & YM: permanently coupled; NS: network cascade \\
$\Phi$ & $\Phi_c$ & $\Phi_c$ & Both critical --- matched \\
\textbf{$\Omega$} & \textbf{$\Omega_Z$} & \textbf{$\Omega_0$} & \textbf{YM: protected; NS: unprotected} \\
P & P\_pm\_sym & P\_asym & YM: symmetric; NS: asymmetric \\
F & F\_hbar & F\_ell & YM: quantum; NS: classical \\
\bottomrule
\end{tabularx}
\end{table}

Both theories operate at $\Phi_c$. Yang-Mills has a mass gap; Navier-Stokes has the blowup question. The structural difference that distinguishes them is $\Omega_Z$ vs $\Omega_0$. The Yang-Mills gap is the physical realization of $\Omega_Z$ topological protection. Navier-Stokes has no equivalent. The grammar predicts: wherever $\Omega_Z$ is absent and P\_asym + $\Phi_c$ are present, there is no structural barrier analogous to the mass gap.

\textbf{S.3} --- $\Gamma_\text{and}$ as Richardson cascade encoding: The Navier-Stokes nonlinear term u$\cdot$$\nabla$u couples all modes simultaneously --- this is the $\Gamma_\text{and}$ signature. You cannot solve Navier-Stokes scale by scale because $\Gamma_\text{and}$ requires all scales to be satisfied at once. The Richardson energy cascade (large eddies $\to$ medium eddies $\to$ small eddies $\to$ dissipation) is the physical realization of $\Gamma_\text{and}$ conjunctive cross-scale coupling. Turbulence is not disordered --- it is $\Gamma_\text{and}$ structure made visible.

\textbf{S.4} --- The decisive empirical question: Does blowup occur in 3D NS? If yes, it confirms the $\Omega_0$ structural argument (T.3--T.4). If no (global smooth solutions exist), it confirms the $\Phi$-cascade argument (T.5) and implies that $\Phi$ dynamics can dominate $\Omega$ protection even at $\Omega_0$. The Millennium problem is, structurally, a question about which primitive wins when $\Phi_\text{cascade}$ and $\Omega_0$ are in competition.

\textbf{What Diaphorics does NOT say:} It does not compute blowup rates vs cascade rates. That quantitative comparison is the remaining mathematical work.

\begin{center}
\rule{0.5\textwidth}{0.4pt}
\end{center}

\subsubsection{Ontics Plane --- Ontological Implications}

\textbf{O.1} --- $\Omega_0$ means NS has no mass-gap analog: {[}structural{]} If the Yang-Mills gap is the $\Omega_Z$ effect --- topological protection manifesting as a minimum energy cost --- then $\Omega_0$ in NS means there is no analogous protection. The grammar predicts blowup is \emph{structurally unprotected}, not merely mathematically unproven. This is a stronger statement than ''''we do not know.'''' The grammar says: the protection mechanism that exists in YM does not exist in NS.

\textbf{O.2} --- The framework is internally divided, and that is informative: {[}structural{]} The two competing arguments ($\Phi$-cascade for smoothness, $\Omega_0$ for blowup) are not a failure of the grammar --- they identify exactly where the mathematical uncertainty lives. The question ''''does $\Phi_\text{super}$ decay outpace singularity formation?'''' is the quantitative gap the framework cannot fill. The primitive structure has narrowed the question from ''''we do not know'''' to ''''we know which two primitives are in competition and which measurement would resolve it.''''

\textbf{O.3} --- Turbulence is $\Gamma_\text{and}$ structure, not disorder: {[}ontological claim, strong{]} The turbulent cascade is not the breakdown of fluid order --- it is the physical realization of $\Gamma_\text{and}$ conjunctive cross-scale coupling operating at $\Phi_c$. Turbulence looks chaotic from D\_wedge (local, scale-specific observation); from D\_infty (the full solution space), it is a structured $\Gamma_\text{and}$ operation propagating constraint through all scales simultaneously. The ''''chaos'''' of turbulence is a G-scope misidentification.

\textbf{What Ontics does NOT say:} It does not settle the existence question (Diaphorics), and it does not claim the blowup argument is proven (that remains a Topics gap --- the rate comparison).

\begin{center}
\rule{0.5\textwidth}{0.4pt}
\end{center}

\subsection{The Separation in Practice: What You Lose by Conflating Planes}

\textbf{Conflating Topics and Diaphorics} produces unfalsifiable grammar claims. Example: ''''T\_bowtie requires minimum energy'''' --- is this true because of the axioms (Topics, testable by checking axiom consistency) or because QCD works that way (Diaphorics, testable by encoding QCD differently)? If you conflate, you cannot tell which evidence would falsify the claim.

\textbf{Conflating Diaphorics and Ontics} produces overreach. Example: ''''The cosmological constant is derivable from first principles'''' --- is this a catalog calculation (Diaphorics, requires knowing E\_Planck and $E_{\Lambda}$ as inputs) or an ontological claim (Ontics, ''''physical constants are G-scope readings'''')? The Diaphorics calculation is currently circular. The Ontics implication is not --- it would hold even if the numbers were slightly different.

\textbf{Conflating Topics and Ontics} produces the worst failure: deriving philosophical conclusions directly from grammar axioms without empirical grounding. Example: ''''Consciousness is structurally guaranteed at $\Phi_c$'''' --- is this a grammar claim (Topics: $\Phi_c$ systems satisfy Axiom 5 self-reference) or an ontological claim (Ontics: self-reference \emph{is} consciousness)? The grammar says self-reference occurs; the grammar cannot say what self-reference \emph{is like} from the inside. The Grammar-Phenomenology Gap {[}ONTO:$\S$IV{]} exists precisely because Topics and Ontics are perpendicular.

\textbf{The rule}: every claim in PRIMITIVE\_THEOREMS.md can be assigned exactly one primary plane. Claims that appear to belong to multiple planes are usually compound --- decompose them.

\begin{center}
\rule{0.5\textwidth}{0.4pt}
\end{center}

\subsection{Template for Future Theorems}

For each theorem in PRIMITIVE\_THEOREMS.md, the three-plane analysis should provide:

\begin{lstlisting}
Topics: [list grammar claims --- derivable from axioms alone]
        [identify which axiom or primitive definition each follows from]
        [state what falsifies each: which axiom must be wrong if this fails]

Diaphorics: [give the primitive encoding of the relevant physical system]
            [cite the empirical data that confirms or could falsify]
            [note any circularities explicitly]

Ontics: [state what the structural result implies about the nature of being/reality]
        [mark each implication: structural (follows necessarily) or interpretive (compatible with)]
        [state explicitly what the result does NOT say ontologically]
\end{lstlisting}

\begin{center}
\rule{0.5\textwidth}{0.4pt}
\end{center}

\emph{This document accompanies PRIMITIVE\_THEOREMS.md and feeds into the canonical document updates:}
\emph{--- Topics claims $\to$ new sections in SYNTHONICON\_TOPICS.md}
\emph{--- Diaphorics claims $\to$ new catalog entries in SYNTHONICON\_DIAPHORICS.md}
\emph{--- Ontics claims $\to$ new sections in SYNTHONICON\_ONTICS.md}


\end{document}
